\chapter*{本次译者补注索引}
\addcontentsline{toc}{chapter}{本次译者补注索引}
本次补注依据 2026 年 11 月 5 日收到的勘误说明编写，编号 1--29 与该说明一一对应。各注列出说明所建议的修正式、适用条件及理由；第 2 项单列证明范围的缺口。作者正式勘误、原有译注与本次补注分别保留。

本版正文保留原书表述，相关位置以“译者补注”说明应如何理解或修正；同一问题再次出现时，另附跳转提示。这里使用章节和公式号定位，并列出本版实际页码。单击条目标题可跳至补注，单击注内的补注编号可返回本索引。

\begin{enumerate}[label=\arabic*.,leftmargin=2.5em,itemsep=7pt]
\item \hypertarget{translator-index-1}{}\hyperlink{translator-note-1}{\csuse{translator-note-title-1}}。
第一篇，命题1.1后的Picard下降论证；第\,\pageref*{translator-note:1}\,页。
\item \hypertarget{translator-index-2}{}\hyperlink{translator-note-2}{\csuse{translator-note-title-2}}。
第一篇，定理6.4及其证明；第\,\pageref*{translator-note:2}\,页。
\item \hypertarget{translator-index-3}{}\hyperlink{translator-note-3}{\csuse{translator-note-title-3}}。
第一篇，备注8.4.2；第\,\pageref*{translator-note:3}\,页。
\item \hypertarget{translator-index-4}{}\hyperlink{translator-note-4}{\csuse{translator-note-title-4}}。
第一篇，定理9.1证明的第二个交换图及说明；第\,\pageref*{translator-note:4}\,页。
\item \hypertarget{translator-index-5}{}\hyperlink{translator-note-5}{\csuse{translator-note-title-5}}。
第二篇，I §2.2.3(3)；第\,\pageref*{translator-note:5}\,页。
\item \hypertarget{translator-index-6}{}\hyperlink{translator-note-6}{\csuse{translator-note-title-6}}。
第二篇，I §2.2.4；第\,\pageref*{translator-note:6}\,页。
\item \hypertarget{translator-index-7}{}\hyperlink{translator-note-7}{\csuse{translator-note-title-7}}。
第二篇，I 猜想2.2.7(ii)及§2.2.6(4)；第\,\pageref*{translator-note:7}\,页。
\item \hypertarget{translator-index-8}{}\hyperlink{translator-note-8}{\csuse{translator-note-title-8}}。
第二篇，I §2.2.9；第\,\pageref*{translator-note:8}\,页。
\item \hypertarget{translator-index-9}{}\hyperlink{translator-note-9}{\csuse{translator-note-title-9}}。
第二篇，I §3.3.5；第\,\pageref*{translator-note:9}\,页。
\item \hypertarget{translator-index-10}{}\hyperlink{translator-note-10}{\csuse{translator-note-title-10}}。
第二篇，I 引理3.5.5；第\,\pageref*{translator-note:10}\,页。
\item \hypertarget{translator-index-11}{}\hyperlink{translator-note-11}{\csuse{translator-note-title-11}}。
第二篇，II §2.1.4；第\,\pageref*{translator-note:11}\,页。
\item \hypertarget{translator-index-12}{}\hyperlink{translator-note-12}{\csuse{translator-note-title-12}}。
第二篇，II §2.2.5；第\,\pageref*{translator-note:12}\,页。
\item \hypertarget{translator-index-13}{}\hyperlink{translator-note-13}{\csuse{translator-note-title-13}}。
第二篇，II 引理2.2.8及§2.2.10证明末行；第\,\pageref*{translator-note:13}\,页。
\item \hypertarget{translator-index-14}{}\hyperlink{translator-note-14}{\csuse{translator-note-title-14}}。
第二篇，II §2.3.1；第\,\pageref*{translator-note:14}\,页。
\item \hypertarget{translator-index-15}{}\hyperlink{translator-note-15}{\csuse{translator-note-title-15}}。
第二篇，III §1.2.5的de Rham同构；第\,\pageref*{translator-note:15}\,页。
\item \hypertarget{translator-index-16}{}\hyperlink{translator-note-16}{\csuse{translator-note-title-16}}。
第二篇，III §1.2.5的辅助映射$a'$；第\,\pageref*{translator-note:16}\,页。
\item \hypertarget{translator-index-17}{}\hyperlink{translator-note-17}{\csuse{translator-note-title-17}}。
第二篇，III 定理1.2.6证明中的局部计算；第\,\pageref*{translator-note:17}\,页。
\item \hypertarget{translator-index-18}{}\hyperlink{translator-note-18}{\csuse{translator-note-title-18}}。
第二篇，III 定理1.2.6证明中的系数代数；第\,\pageref*{translator-note:18}\,页。
\item \hypertarget{translator-index-19}{}\hyperlink{translator-note-19}{\csuse{translator-note-title-19}}。
第三篇，定义2.3(d)及定理2.4；第\,\pageref*{translator-note:19}\,页。
\item \hypertarget{translator-index-20}{}\hyperlink{translator-note-20}{\csuse{translator-note-title-20}}。
第三篇，§7第一阶段的Riemann--Roch下界；第\,\pageref*{translator-note:20}\,页。
\item \hypertarget{translator-index-21}{}\hyperlink{translator-note-21}{\csuse{translator-note-title-21}}。
第三篇，§7第二阶段的单项式范围；第\,\pageref*{translator-note:21}\,页。
\item \hypertarget{translator-index-22}{}\hyperlink{translator-note-22}{\csuse{translator-note-title-22}}。
第三篇，引理8.1；第\,\pageref*{translator-note:22}\,页。
\item \hypertarget{translator-index-23}{}\hyperlink{translator-note-23}{\csuse{translator-note-title-23}}。
第三篇，Faltings复形(9.3)；第\,\pageref*{translator-note:23}\,页。
\item \hypertarget{translator-index-24}{}\hyperlink{translator-note-24}{\csuse{translator-note-title-24}}。
第三篇，式(13.10)及推论18.3证明；第\,\pageref*{translator-note:24}\,页。
\item \hypertarget{translator-index-25}{}\hyperlink{translator-note-25}{\csuse{translator-note-title-25}}。
第三篇，命题14.1证明中的$\Gamma_1$；第\,\pageref*{translator-note:25}\,页。
\item \hypertarget{translator-index-26}{}\hyperlink{translator-note-26}{\csuse{translator-note-title-26}}。
第三篇，引理15.3及其前置设定；第\,\pageref*{translator-note:26}\,页。
\item \hypertarget{translator-index-27}{}\hyperlink{translator-note-27}{\csuse{translator-note-title-27}}。
第三篇，引理15.8；第\,\pageref*{translator-note:27}\,页。
\item \hypertarget{translator-index-28}{}\hyperlink{translator-note-28}{\csuse{translator-note-title-28}}。
第三篇，定理18.1、推论18.3及式(18.2)、(18.4)；第\,\pageref*{translator-note:28}\,页。
\item \hypertarget{translator-index-29}{}\hyperlink{translator-note-29}{\csuse{translator-note-title-29}}。
第三篇，范数估计(18.5)及后续对数估计；第\,\pageref*{translator-note:29}\,页。
\end{enumerate}
